\scanpage{212}

\deleted{pour simplifier, $S$ \inserted{loc.} noeth, donc les $S'$ \deleted{\unsure{tous}} aussi (NB ils sont l. de p.f. sur $S$) \unsure{du lieu|du leur|du livre|dans l'énoncé} \unsure{nilradical|nilradicaux|nilpotents} de \unsure{ce quotient|ce qui précède|(sur $S$)}, on trouve}

On voit que, une connexion à courbure nulle sur un Mod. inversible sur $X_{S'}/S'$ est à courbure \uline{nilpotente} ssi, en restreignant sur $S'_{\mathrm{réd}}$, on trouve connexion compatible aux puissances $p^{\text{ièmes}}$. Donc on trouve que $P^{cin}_{X/S}$ est l'image inverse dans $P^{ci}_{X/S}$ du sous faisceau $\overline{f_*^{\mathrm{fl}}(\Omega^{1\otimes p}_{X/S})}\subset f_*^{\mathrm{fl}}(\Omega^{1\otimes p}_{\unsure{X|x}/S})$, qui est le \deleted{\unsure{sous}} ``grpe $\pm$ formel associé'' \ldots. On trouve ainsi

\[
0\longrightarrow P^{cip}\longrightarrow P^{ci}_{X/S}\longrightarrow f_*^{\mathrm{fl}}(\Omega^{1\otimes p}_{X/S})/\overline{f_*^{\mathrm{fl}}(\Omega^{1\otimes p}_{X/S})}.\tag*{(4.1)}
\]

\deleted{\unsure{4}} et (3.1) devient

\diagram{\begin{tikzcd}[ampersand replacement=\&,column sep=small,row sep=tiny]
\llap{(4.2)\qquad}0\arrow[r]\&P^{cip}_{X/S}\arrow[r]\&P^{cin}_{X/S}\arrow[r]\&\overline{f_*^{\mathrm{fl}}(\Omega^{1\otimes p}_{X/S})},\\
\& \underline{\mathrm{Pic}}_{X/S}^{p/S}\arrow[u,phantom,"\rotatebox{90}{$\simeq$}" description]
\end{tikzcd}}

qui explicite $P^{cin}$ comme extension d'un ss-grpe d'un grpe $\pm$formel par \deleted{\unsure{un}} le grpe $P^{cip}$, qui est bien compris \ldots

\deleted{On peut aussi se proposer d'étudier}

On conclut

\diagram{\begin{tikzcd}[ampersand replacement=\&,column sep=small,row sep=tiny]
\llap{(4.3)\qquad}0\arrow[r]\&P^{cip\,\unsure{0|o}}_{X/S}\arrow[r,dash]\&P^{cin\,0}_{X/S}\arrow[r]\&\overline{f_*^{\mathrm{fl}}(\Omega^{1\otimes p}_{X/S})}.\\
\&\underline{\mathrm{Pic}}_{X/S}^{0\,p/S}\arrow[u,phantom,"\rotatebox{90}{$\simeq$}" description]
\end{tikzcd}}

\deleted{On peut aussi aller en sens inverse, et}

étudier le noyau $U$ \deleted{\unsure{de}} \inserted{\unsure{dans|de}} $P^{cin}_{X/S}\to\underline{\mathrm{Pic}}_{X/S}$,

\[
0\longrightarrow U\longrightarrow P^{cin}_{X/S}\longrightarrow\underline{\mathrm{Pic}}_{X/S}\tag*{(4.4)}
\]

c'est le noyau de $\dfrac{f_*^{\mathrm{fl}}(\underline Z^1_{X/S})}{\mathrm{Im}\,f_*^{\mathrm{fl}}(O^*_{X})}\longrightarrow f_*^{\mathrm{fl}}(\Omega^{1\otimes p}_{X/S})/\overline{f_*^{\mathrm{fl}}(\Omega^{1\otimes p}_{X/S})}$, donc il admet une dévissage en

\[
0\longrightarrow\underline{\mathrm{Cart}}_{X/S}\longrightarrow U\longrightarrow\overline{f_*^{\mathrm{fl}}(\Omega^{1\otimes p}_{X/S})}\tag*{(4.5)}
\]

% archive sheet number: 211
